\addcontentsline{toc}{chapter}{Conclusion générale}
\chapter*{Conclusion générale}

Au terme de ces travaux, nous avons développé Numeris229, une solution de RAG souveraine et spécialisée, capable de transformer les interrogations en langage naturel des citoyens et des professionnels en recherches juridiques précises au sein du Code du Numérique en République du Bénin. L’objectif principal était de démocratiser, sécuriser et fiabiliser l’accès aux dispositions de la Loi N° 2017-20 modifiée, souvent perçues comme denses et complexes. Grâce à une approche hybride associant des représentations sémantiques francophones, un filtrage lexical exact et un ré-ordonnancement par cross-encoder, le système parvient à éliminer les hallucinations des modèles de langage tout en guidant les usagers vers les articles officiels régissant rigoureusement leur situation. Les évaluations conduites ont démontré qu'une telle chaîne d'inférence locale offre une précision documentaire supérieure aux agents conversationnels généralistes. Cependant, au-delà de ces résultats encourageants, des défis subsistent, notamment l’intégration des décrets d'application, l'extension vers des canaux de messagerie instantanée comme WhatsApp et la traduction vocale en langues nationales. En somme, ce projet prouve que l’intelligence artificielle, lorsqu’elle est rigoureusement conçue et ancrée dans le cadre légal local, peut constituer un levier d'inclusion juridique puissant, fiable et déterminant pour l'appropriation du droit numérique par l'ensemble de la société béninoise.